\thispagestyle{empty} % Asegura que esta página no tenga números ni encabezados/pies de página
\begin{titlepage}
	
	% 1. Logo/Incio de imagen (alineado a la derecha en la parte superior)
	\begin{flushright}
		\vspace*{0.5cm} % Espacio desde el borde superior
		
		% Coloca \includegraphics directamente aquí, sin el entorno figure
		\includegraphics[width=0.2\linewidth]{images/Logo}
		
		\vspace{0.2cm} % Espacio entre la imagen y el texto
	\end{flushright}
	
	\vfill % Empuja el contenido restante hacia abajo
	
	% 2. Contenido principal (alineado a la izquierda, cerca de la parte inferior)
	% Franja gris con el título
	\colorbox{franjagris}{
		\begin{minipage}{\textwidth}
			\raggedleft % CAMBIO AQUÍ: Alinea el texto a la derecha
			\vspace{0.8cm} % Espacio superior dentro de la caja
			\Large \textbf{WaveLink32:} \\
			\vspace{0.2cm}
			\Huge \textbf{Creación de un carro dirigido por comandos de voz} \\
			\vspace{0.8cm} % Espacio inferior dentro de la caja
		\end{minipage}
	}
	
	\vspace{0.5cm} % Espacio entre la franja y el texto del pie
	
	\raggedleft 
	Documentación sobre la elaboración de un carro controlado por comandos de voz procesados por Matlab. La lógica de programación fue creada por medio de Simulink y el EmbeddedCoder.
\end{titlepage}